% Appendix H: comparison with other CT collections (promised to reviewer Fbco in the author
% response of 27 Jul 2026). Integrated 2026-09-27 (third pass) from the verified draft
% scratchpad/neurips_promises/drafts/app_comparison.tex; the licence line now also gives ours
% per host, and note g points to the total_v3 slot.
%
% Needs only what neurips_2026.tex already loads (booktabs; amssymb for \checkmark) and its
% macros \dataset, \TS, \VerSe, \VERIDAH. Its bib entries are merged into bib/references.bib.
% Every cell and its source: drafts/app_comparison_NOTES.md.
% The CTSpinoPelvic1K row is the April 2026 release (1,153 records, 802 patients).
% Label tab:dataset_comparison (tab:comparison is taken by the baselines table).
\section{Comparison with other CT collections}
\label{app:comparison}

Table~\ref{tab:dataset_comparison} compares \dataset{} with the CT collections that label
anatomy at the lumbosacral junction; each cell was checked against the collection's own
paper or data page. \VerSe{} comes closest in intent: it labels L6 and T13, lists anatomical
variants per subject and grades Castellvi II--IV, but it segments neither the sacrum nor a
transitional vertebra partially fused to it, so its type~III and IV cases cannot be scored
at L5/S1. CTSpine1K labels L6 but no pelvis, and its nnU-Net baseline already scored L6 at
a Dice of 0 on all three of its test sets, with shifted predictions on lumbarization and
sacralization cases (its Table~2 and Fig.~4)~\citep{Deng2021CTSpine1K:Tomography}, which is
consistent with the unmerged control of \S\ref{sec:unmerged}. Among these collections, only
\dataset{} puts numbered lumbar levels including L6, the sacrum and both hip bones on the
same patients, with a per-patient LSTV label and, for its 33 LSTV patients, a Castellvi grade.

\begin{table}[h]
\caption{CT collections labelling the lumbosacral junction. \checkmark{}
present; (\checkmark) partial; -- absent; letters refer to the notes. Patients in
parentheses where reported and different from scans. Pelvis: S sacrum, H hip bones,
F femora. Unified: spine and pelvis in the same patients. Adapted from Table~I of the
dataset article~\citep{Schwing2026CTSpinoPelvic1KArxiv}, which describes a later release.}
\label{tab:dataset_comparison}
\centering
\small
\setlength{\tabcolsep}{3pt}
\begin{tabular}{@{}lrlcccccc@{}}
\toprule
\bfseries Collection & \bfseries Scans & \bfseries Levels & \bfseries Pelvis
  & \bfseries Unified & \bfseries LSTV & \bfseries Six-way & \bfseries Castellvi
  & \bfseries Public \\
 & (patients) & labelled & & cohort & labels & subtype & grade & data \\
\midrule
\VerSe{}
  & 374 (355) & C1--L6, T13\textsuperscript{a} & --
  & -- & (\checkmark)\textsuperscript{b} & -- & (\checkmark)\textsuperscript{c} & \checkmark \\
CTSpine1K
  & 1{,}005 & C1--L6\textsuperscript{d} & --
  & -- & (\checkmark)\textsuperscript{d} & -- & -- & \checkmark \\
CTPelvic1K
  & 1{,}184 & lumbar, 1 class\textsuperscript{e} & S, H
  & (\checkmark)\textsuperscript{e} & (\checkmark)\textsuperscript{f} & -- & -- & \checkmark \\
\TS{}
  & 1{,}204 & C1--L5\textsuperscript{g} & S, H, F
  & (\checkmark)\textsuperscript{g} & -- & -- & -- & \checkmark \\
\VERIDAH{} CT cohort
  & 1{,}536 & C1--L6, T13\textsuperscript{h} & --\textsuperscript{h}
  & -- & (\checkmark)\textsuperscript{h} & -- & --\textsuperscript{h} & -- \\
\dataset{}
  & 1{,}153 (802)\textsuperscript{i} & L1--L6 & S, H
  & \checkmark\textsuperscript{i} & \checkmark\textsuperscript{j} & \checkmark\textsuperscript{j}
  & \checkmark\textsuperscript{k} & \checkmark\textsuperscript{i} \\
\bottomrule
\end{tabular}

\smallskip
\footnotesize
\raggedright
\textsuperscript{a}Identifiers 1--25 (C1--L6) and 28 (T13); sacrum and coccyx not
labelled~\citep{VerSeRepo}; 374 scans of 355 patients~\citep{Sekuboyina2021VerSe:Images}.
\textsuperscript{b}VerSe'20 was assembled partly from reports of anatomical variants and
lists them per subject; 85 of its cases have six lumbar vertebrae~\citep{Liebl2021VerSeData}.
\textsuperscript{c}Castellvi II--IV graded in VerSe'20 (IIa 35, IIb 25, IIIa 14, IIIb 28,
IV 7), types 0 and I pooled; the sacrum and vertebrae partially fused to it (III, IV) are
not segmented~\citep{Liebl2021VerSeData}.
\textsuperscript{d}Identifiers 1--25 (C1--L6); 784 of the 1{,}005 volumes are COLONOG.
Its Table~1 counts 17 lumbarization and 11 sacralization cases, ungraded; the authors list
their IDs with the data~\citep{Deng2021CTSpine1K:Tomography}.
\textsuperscript{e}Four classes: lumbar spine (one class), sacrum, left and right hip, on
the same scan, cropped to the pelvic region; patient count not
reported~\citep{Liu2021DeepModels}.
\textsuperscript{f}Undocumented tags in ten of the 713 COLONOG mask filenames on Zenodo:
sacralization 7, semisacralization 2, hard\_sacralization 1~\citep{Liu2020CTPelvic1KData,Schwing2026CTSpinoPelvic1KArxiv}.
\textsuperscript{g}The cited paper's dataset (v1, 104 classes)~\citep{Wasserthal2023TotalSegmentator:Images,Wasserthal2022TSDatasetV1};
v2.0.1 has 1{,}228 scans and 117 classes and adds
S1~\citep{Wasserthal2023TSDatasetV2,TotalSegmentatorCode}; neither has L6 (the software's
later \texttt{total\_v3} task has an L6 class; Appendix~\ref{app:tsv3}). Every
structure in the field of view is labelled; how many scans hold both lumbar spine and
pelvis is not reported.
\textsuperscript{h}In-house cohort, patient count not reported; reference labels include
T13 and L6, merged into T12 and L5 for training; enumeration anomalies are labelled by
count, and a Castellvi transitional vertebra is labelled sacralized lumbar whatever its type;
lumbar and sacral vertebrae are distinguished, but no sacral or pelvic segmentation is
described; code and weights public, no image release
described~\citep{Moller2026VERIDAH:Sequences}.
\textsuperscript{i}Records of the April 2026 release used here (Appendix~\ref{app:release}).
Spine and pelvis for 693 of 802 patients: on one series for 342, on two acquisitions of the
same patient (two records) for 351; 89 spine-only, 20 pelvis-only. Public as v11 (the same
802 patients, one record each); the April records and evaluated split are not yet at a
cited location.
\textsuperscript{j}For all 802 patients (33 LSTV), from the source masks' lumbar counts and
filename tags, not image reads (\S\ref{sec:dataset}); the six subtypes are listed per
patient in the v11 manifest (\texttt{lstv\_phenotype}).
\textsuperscript{k}Types I--IV for all 33 LSTV patients: a two-resident consensus completed
in September 2026, after the April release, and distributed from Zenodo v8 onward
(\S\ref{sec:castellvi}, Appendix~\ref{app:castellvi}); normal patients not screened for
type I.
Licences: \VerSe{} CC BY-SA 4.0; CTSpine1K CC BY-NC-SA 4.0; CTPelvic1K CC BY-NC-SA (paper)
or CC BY 4.0 (Zenodo record); \TS{} data CC BY 4.0; \dataset{} CC BY-NC-SA 4.0 on Zenodo, the
archive of record, and CC BY-NC 4.0 on the dataset card of its HuggingFace mirror (checked
27 September 2026).
\end{table}

\paragraph{How the benchmark will be used.}
Its first use is as a stratified test: a CT segmenter or vertebra labeller is scored per
LSTV subtype with junction-DSC and voxel confusion at L5/S1,
each stratum reported with its $n$ (\S\ref{sec:protocol}), which is what a model claiming
readiness for surgical planning should show. Adding an L6 class did not remove the L5/L6
collision at this sample size (\S\ref{sec:unmerged}), and the signal that separates a last
lumbar vertebra from its neighbour is a count that one patch does not reliably contain, so
the benchmark is meant as the test set for counting-aware and two-stage ``segment, then
number'' designs, such as a segmenter feeding \VERIDAH{}'s sequence predictor
(\S\ref{sec:discussion}). Its LSTV strata are small (33 patients, 7 with an LSTV-tagged
held-out record; Wilson intervals for strata of 3--4 cases start at 44--51\%,
\S\ref{sec:limits-test-set}), so it exposes a failure mode and its direction rather than
ranking methods that differ by a few points.
